\chapter{Trabajo y energía mecánica}\label{ch:g12:work-and-energy}

Echa una canica en una ensaladera: se lanza hacia abajo, se pasa del fondo, sube hasta
su altura de partida, duda y da media vuelta --- y cada inversión puede predecirse a
partir de una curva (el perfil del cuenco) y una recta horizontal (la \omterm{def:g10:energy-conservation:energy}{energía} de la
canica), sin resolver ninguna ecuación del movimiento. El curso pasado puso precio a
las \omterm{def:g10:inertia:force}{fuerzas} sobre trayectos rectilíneos (\cref{ch:g11:work-of-force}) y equilibró la
\omterm{def:g11:mechanical-energy:kinetic}{energía cinética} con la potencial (\cref{ch:g11:mechanical-energy}); este año la
contabilidad llega a todas partes: caminos curvos, resortes, paisajes de un vistazo.

\section{El trabajo a lo largo de un camino cualquiera}

\begin{definition}[Trabajo a lo largo de un camino]\label{def:g12:work-and-energy:path}
Sea una \omterm{def:g10:inertia:force}{fuerza} $\vect F$ que actúa sobre un cuerpo que se mueve por un camino
cualquiera de $A$ a $B$: corta el camino en desplazamientos $\vect{\dd\ell}$ tan
cortos que cada uno sea recto y $\vect F$ apenas cambie a lo largo de él. Cada paso se
gana el pequeño \omterm{def:g11:work-of-force:work}{trabajo} $\scal{F}{\dd\ell} = F\,\dd\ell\cos\theta$
(\cref{ch:g11:work-of-force}), y el \emph{\omterm{def:g11:work-of-force:work}{trabajo} a lo largo del camino}%
\index{trabajo (de una fuerza)!a lo largo de un camino} es su suma,
$W_{AB}(\vect F) = \sum \scal{F}{\dd\ell}$ --- que sobre un camino recto y con \omterm{def:g10:inertia:force}{fuerza}
constante se reduce al $F \times AB \times \cos\theta$ del curso pasado.
\end{definition}

\begin{omfigure}
\begin{tikzpicture}[font=\small]
  \draw[omExo!60, thick] plot [smooth, tension=0.8] coordinates {(0,0.3) (1.2,1.15) (2.6,1.3) (4.0,0.7) (5.3,1.05)};
  \draw[omMeth, thick] (0,0.3) -- (1.2,1.15) -- (2.6,1.3) -- (4.0,0.7) -- (5.3,1.05);
  \foreach \p in {(0,0.3),(1.2,1.15),(2.6,1.3),(4.0,0.7),(5.3,1.05)} \fill \p circle (1.3pt);
  \node[below] at (0,0.3) {$A$}; \node[below right] at (5.3,1.05) {$B$};
  \draw[-Stealth, omThm, very thick] (2.6,1.3) -- (3.44,0.94) node[below left=-2pt] {$\vect{\dd\ell}$};
  \draw[-Stealth, omDef, very thick] (2.6,1.3) -- ++(22:1.5) node[above] {$\vect F$};
  \draw (2.6,1.3) ++(-23:0.75) arc[start angle=-23, end angle=22, radius=0.75];
  \node at ($(2.6,1.3)+(0:1.0)$) {$\theta$};
\end{tikzpicture}

{\small Un camino curvo es una línea quebrada de muchos pasos cortos: cada
$\vect{\dd\ell}$ se gana $F\,\dd\ell\cos\theta$; el \omterm{def:g11:work-of-force:work}{trabajo} es la suma.}
\end{omfigure}

\begin{proposition}[Trabajo del peso, por cualquier camino]\label{prop:g12:work-and-energy:weight}
A lo largo de \emph{cualquier} camino de $A$ (altitud $z_A$) a $B$ (altitud $z_B$), el
\omterm{def:g10:universal-gravitation:weight}{peso} de una masa $m$ realiza el \omterm{def:g11:work-of-force:work}{trabajo} $W_{AB}(\vect P) = mg\,(z_A - z_B)$: entra el
desnivel, nunca el recorrido.
\end{proposition}

\begin{proof}
Cada paso corto aporta $mg \times (\text{desnivel pequeño})$ (el cálculo del curso
pasado sobre segmentos rectos); los desniveles se telescopan hasta $z_A - z_B$. El
caso curvo, admitido entonces, queda entregado.
\end{proof}

\begin{definition}[Fuerzas conservativas y no conservativas]\label{def:g12:work-and-energy:conservative}
Una \omterm{def:g10:inertia:force}{fuerza} es \emph{conservativa}\index{fuerza conservativa} si su \omterm{def:g11:work-of-force:work}{trabajo} entre dos
puntos es el mismo por cualquier camino (nulo en todo trayecto de ida y vuelta), y
\emph{no conservativa}\index{fuerza no conservativa} en caso contrario. El \omterm{def:g10:universal-gravitation:weight}{peso} y la
\omterm{def:g10:inertia:force}{fuerza} de un resorte (sección siguiente) son conservativos. El \omterm{def:g10:inertia:inventory}{rozamiento} de
deslizamiento no lo es: opuesto al movimiento, realiza $W = -fL$ sobre un camino de
longitud $L$ --- un peaje por \omterm{def:g10:orders-of-magnitude:unit}{metro} que no se devuelve jamás.
\end{definition}

\begin{example}[Dos senderos hasta el refugio]\label{ex:g12:work-and-energy:trails}
Un excursionista de \qty{75}{kg} gana \qty{300}{m} por un sendero empinado de
\qty{800}{m} o por \qty{2.4}{km} de zigzags. El \omterm{def:g10:universal-gravitation:weight}{peso} le cobra
$-75 \times 9.81 \times 300 \approx \qty{-2.2e5}{J}$ por los dos; una \omterm{def:g10:inertia:force}{fuerza} de
\omterm{def:g10:inertia:inventory}{rozamiento} de \qty{30}{N}, \qty{-2.4e4}{J} por uno y \qty{-7.2e4}{J} por el otro. El
\omterm{def:g10:universal-gravitation:weight}{peso} devuelve lo suyo en la bajada; la factura del \omterm{def:g10:inertia:inventory}{rozamiento} se ha ido en calor, en
ambos sentidos.
\end{example}

\section{La energía potencial}

\begin{definition}[Energía potencial]\label{def:g12:work-and-energy:ep}
A cada \omterm{def:g12:work-and-energy:conservative}{fuerza conservativa} se le asocia una \emph{\omterm{def:g10:energy-conservation:energy}{energía} potencial}%
\index{energía potencial!a partir del trabajo conservativo} $E_p$, un número que
depende solo de la posición y que se define mediante $W_{AB} = E_p(A) - E_p(B)$: el
\omterm{def:g11:work-of-force:work}{trabajo} entregado es \omterm{def:g10:energy-conservation:energy}{energía} potencial gastada. Fijar $E_p = 0$ en un punto de
referencia elegido la fija en todas partes; la elección desplaza todos los valores en
una constante que se cancela en cualquier diferencia. Para el \omterm{def:g10:universal-gravitation:weight}{peso} esto devuelve el
$E_p = mgz$ del curso pasado, válido ahora por cualquier camino. El \omterm{def:g10:inertia:inventory}{rozamiento} no
puede tener $E_p$: una factura que depende del recorrido no es una diferencia de
números asociados a los extremos --- solo las \omterm{def:g12:work-and-energy:conservative}{fuerzas conservativas} almacenan de forma
recuperable, que es exactamente lo que el nombre conserva.
\end{definition}

\begin{proposition}[Energía potencial elástica]\label{prop:g12:work-and-energy:elastic}
Un resorte de \omterm{def:g12:oscillators-and-time:spring}{constante elástica} $k$ (\unit{N/m}) estirado o comprimido una longitud
$x$ tira de vuelta con la \omterm{def:g10:inertia:force}{fuerza} $F = kx$ (\cref{ch:g12:oscillators-and-time}); el
\omterm{def:g11:work-of-force:work}{trabajo} que entrega al volver al reposo --- su \emph{\omterm{def:g10:energy-conservation:energy}{energía} potencial elástica}%
\index{energía potencial!elástica} --- vale $E_p = \tfrac12\, k x^2$.
\end{proposition}

\begin{proof}
La \omterm{def:g12:oscillators-and-time:pendulum}{fuerza recuperadora} actúa a lo largo del movimiento, así que $\sum F\,\dd\ell$ es
el \emph{área bajo la gráfica $F$--$x$}: un triángulo de base $x$ y altura $kx$, de
área $\tfrac12 x \times kx$. La \omterm{def:g10:inertia:force}{fuerza} en cada elongación es la misma a la ida que a
la vuelta: \omterm{def:g11:work-of-force:work}{trabajo} fijado por los extremos, \omterm{def:g10:inertia:force}{fuerza} del resorte \omterm{def:g12:work-and-energy:conservative}{conservativa} y $E_p$
bien definida.
\end{proof}

\begin{omfigure}
\begin{tikzpicture}
\begin{axis}[omaxis, width=8.4cm, height=4.8cm, xmin=0, xmax=0.095,
    ymin=0, ymax=21,
    xlabel={elongación $x$ (\unit{m})}, ylabel={fuerza $F$ (\unit{N})}]
  \addplot[draw=none, fill=omDef!18, domain=0:0.06, samples=2] {200*x} \closedcycle;
  \addplot[omThm, thick, domain=0:0.09, samples=2] {200*x};
  \draw[dashed, omExo] (axis cs:0,12) -- (axis cs:0.06,12) -- (axis cs:0.06,0);
  \node[font=\small, anchor=west] at (axis cs:0.061,6.5) {$kx$};
  \node[font=\small] at (axis cs:0.044,3.2) {$\tfrac12 kx^2$};
  \node[font=\small, anchor=south, omThm] at (axis cs:0.082,17) {$F = kx$};
\end{axis}
\end{tikzpicture}

{\small La \omterm{def:g10:inertia:force}{fuerza} del resorte crece con la elongación; la \omterm{def:g10:energy-conservation:energy}{energía} almacenada es el
área bajo la recta: un triángulo, $\tfrac12 \times x \times kx$.}
\end{omfigure}

\section{Los dos teoremas de la energía}

\begin{theorem}[Teorema de la energía cinética]\label{thm:g12:work-and-energy:ketheorem}\index{teorema de la energía cinética}
Para cualquier movimiento de un cuerpo de masa $m$ de $A$ a $B$ --- curvo, con bucles,
tridimensional ---
\[
\Delta E_k = E_k(B) - E_k(A) = \sum W_{AB}(\vect F),
\]
donde la suma recorre todas las \omterm{def:g10:inertia:force}{fuerzas} que actúan sobre el cuerpo.
\end{theorem}

\begin{proof}
Como $v^2 = \scal{v}{v}$, la regla del producto de las matemáticas de este curso da la
derivada de $v^2$ como $2\,\scal{a}{v}$, luego
$\dd E_k/\dd t = m\,\scal{a}{v} = \scal{F_{\text{tot}}}{v}$ por la \omterm{thm:g12:newtons-laws:second}{segunda ley de
Newton} (\cref{ch:g12:newtons-laws}). En un $\dd t$ corto el cuerpo se desplaza
$\vect{\dd\ell} = \vect v\,\dd t$: $E_k$ gana el pequeño \omterm{def:g11:work-of-force:work}{trabajo} de la
\cref{def:g12:work-and-energy:path}; sumar los pasos suma los \omterm{def:g11:work-of-force:work}{trabajos} --- el
enunciado admitido del curso pasado, ganado por completo.
\end{proof}

\begin{proposition}[Potencia instantánea]\label{prop:g12:work-and-energy:power}
En cada instante, una \omterm{def:g10:inertia:force}{fuerza} $\vect F$ sobre un cuerpo de velocidad $\vect v$ entrega
la \omterm{def:g11:work-of-force:power}{potencia} $P = \scal{F}{v} = Fv\cos\theta$, y $\dd E_k/\dd t$ es la suma de las
\omterm{def:g11:work-of-force:power}{potencias} de todas las \omterm{def:g10:inertia:force}{fuerzas}: la fórmula del curso pasado a velocidad constante,
exacta ahora en cada instante de cualquier movimiento.
\end{proposition}

\begin{proof}
Se lee en la demostración anterior: cada \omterm{def:g10:inertia:force}{fuerza} aporta a $E_k$
$\scal{F}{\dd\ell} = \scal{F}{v}\,\dd t$ durante $\dd t$.
\end{proof}

\begin{theorem}[Teorema de la energía mecánica]\label{thm:g12:work-and-energy:emtheorem}\index{teorema de la energía mecánica}
Sea $E_m = E_k + E_p$, donde $E_p$ reúne las \omterm{def:g10:energy-conservation:energy}{energías} potenciales de todas las \omterm{def:g12:work-and-energy:conservative}{fuerzas
conservativas} en juego. Entonces, a lo largo de cualquier camino,
\[
\Delta E_m = W_{AB}(\text{fuerzas no conservativas}),
\]
y $E_m$ se conserva siempre que las \omterm{def:g12:work-and-energy:conservative}{fuerzas no conservativas} no realicen \omterm{def:g11:work-of-force:work}{trabajo}.
\end{theorem}

\begin{proof}
Desglosemos el \omterm{thm:g11:mechanical-energy:ketheorem}{teorema de la energía cinética}:
$\Delta E_k = W_{\text{cons}} + W_{\text{nc}} = -\Delta E_p + W_{\text{nc}}$; pásese
$\Delta E_p$ al otro lado.
\end{proof}

\begin{example}[El péndulo, esta vez honradamente]\label{ex:g12:work-and-energy:pendulum}
Una masa pendular colgada de un hilo de longitud $L = \qty{2.5}{m}$ se suelta en
reposo a $\qty{45}{\degree}$ de la vertical. La tensión, perpendicular al movimiento,
no tiene \omterm{def:g11:work-of-force:power}{potencia}; el \omterm{def:g10:universal-gravitation:weight}{peso} es conservativo: $E_m$ se conserva \emph{sobre el arco
curvo}, cosa que el curso pasado solo podía comprobarse. La masa parte $h = L(1 -
\cos\qty{45}{\degree}) = \qty{0.73}{m}$ más arriba, así que abajo
$v = \sqrt{2gh} \approx \qty{3.8}{m/s}$.
\end{example}

\begin{method}[Auditoría energética de un movimiento cualquiera]\label{met:g12:work-and-energy:audit}
\begin{enumerate}
  \item Elige el sistema y enumera las \omterm{def:g10:inertia:force}{fuerzas} que actúan sobre él.
  \item Clasifícalas: perpendiculares al movimiento $\to$ no trabajan; \omterm{def:g12:work-and-energy:conservative}{conservativas}
        $\to$ una $E_p$ dentro de $E_m$; el resto $\to$ $W_{\text{nc}}$ (\omterm{def:g10:inertia:inventory}{rozamiento} de
        deslizamiento: $-fL$, con $L$ = longitud del camino).
  \item Escribe $\Delta E_m = W_{\text{nc}}$ entre el punto de datos y el punto de la
        pregunta; resuelve. La \omterm{def:g10:energy-conservation:energy}{energía} responde a ``a qué rapidez, a qué altura'';
        solo el tiempo y la aceleración necesitan las leyes de Newton.
\end{enumerate}
\end{method}

\section{Diagramas de energía}

\begin{definition}[Diagrama de energía, puntos de retorno]\label{def:g12:work-and-energy:diagram}
Para un movimiento a lo largo de un eje $x$ con $E_m$ conservada, el \emph{diagrama de
energía}\index{diagrama de energía} representa la curva $E_p(x)$ y la recta horizontal
$E_m$. Como $E_k = E_m - E_p \geq 0$, el movimiento queda confinado allí donde la
curva está por debajo de la recta; la separación entre ambas es $E_k$, leída
directamente. En un \emph{punto de retorno}\index{punto de retorno}, donde la curva se
encuentra con la recta, $v = 0$ y el movimiento se invierte.
\end{definition}

\begin{omfigure}
\begin{tikzpicture}
\begin{axis}[omaxis, width=9.0cm, height=5.4cm,
    xlabel={$x$}, ylabel={energía}, xtick=\empty, ytick=\empty,
    xmin=-2.25, xmax=2.25, ymin=0, ymax=2.45]
  \addplot[omDef, thick, domain=-2.05:2.05, samples=80] {0.25*x^4 - x^2 + 1.5};
  \addplot[omThm, dashed, thick, domain=-2.25:2.25, samples=2] {1.0};
  \addplot[omProp, dashed, domain=-2.25:2.25, samples=2] {1.7};
  \node[font=\small, omThm, anchor=west] at (axis cs:-2.2,0.86) {$E_m$};
  \node[font=\small, omProp, anchor=west] at (axis cs:-2.2,1.85) {$E_m'$};
  \node[font=\small, omDef, anchor=west] at (axis cs:1.55,2.1) {$E_p(x)$};
  \fill (axis cs:0.765,1.0) circle (1.6pt) node[font=\small, below=1pt] {$x_1$};
  \fill (axis cs:1.848,1.0) circle (1.6pt) node[font=\small, below right=-1pt] {$x_2$};
  \fill[omMeth] (axis cs:1.414,0.5) circle (2.2pt) node[font=\small, below] {estable};
  \fill[omMeth] (axis cs:0,1.5) circle (2.2pt) node[font=\small, above] {inestable};
\end{axis}
\end{tikzpicture}

{\small Un \omterm{def:g12:work-and-energy:diagram}{diagrama de energía}, leído sin resolver nada: con la \omterm{def:g10:energy-conservation:energy}{energía} $E_m$ el
cuerpo oscila entre $x_1$ y $x_2$, y va más rápido donde la separación es mayor; con
$E_m'$ repta por encima de la colina y escapa.}
\end{omfigure}

\begin{proposition}[La fuerza a partir del paisaje; equilibrios]\label{prop:g12:work-and-energy:equilibria}
La \omterm{def:g12:work-and-energy:conservative}{fuerza conservativa} a lo largo de $x$ es menos la pendiente de $E_p$:
$F = -\,\dd E_p/\dd x$ --- apunta \emph{cuesta abajo} en el diagrama y se anula allí
donde la tangente es horizontal: un equilibrio. Un mínimo de $E_p$ es
\emph{estable}\index{equilibrio estable} --- el cuerpo desplazado es empujado de
vuelta, como una bola en un valle; un máximo es \emph{inestable}\index{equilibrio
inestable} --- lo empuja lejos, como una bola en la cima de una colina.
\end{proposition}

\begin{proof}
En un paso pequeño $\dd x$ la \omterm{def:g10:inertia:force}{fuerza} realiza el \omterm{def:g11:work-of-force:work}{trabajo}
$F\,\dd x = E_p(x) - E_p(x + \dd x) = -\dd E_p$; divídase entre $\dd x$. Cuesta abajo,
a un lado y a otro de un valle, se apunta hacia dentro; alrededor de una cima, hacia
fuera.
\end{proof}

\begin{remark}[Oscilaciones y escape]\label{rem:g12:work-and-energy:escape}
Las pequeñas oscilaciones del \cref{ch:g12:oscillators-and-time} viven en el fondo de
los valles de $E_p$. Y lejos de la Tierra la gravedad se debilita: la curva $E_p(r)$
sube cada vez más despacio hacia un techo finito (el área bajo la gráfica de la \omterm{def:g10:inertia:force}{fuerza}
menguante es finita). Un cohete cuya $E_m$ \omterm{def:g12:projectile-motion:range}{alcance} ese techo no encuentra jamás un
\omterm{def:g12:work-and-energy:diagram}{punto de retorno}: \emph{escapa}\index{velocidad de escape} --- alrededor de
\qty{11.2}{km/s} desde el suelo (\cref{exo:g12:work-and-energy:13}); por debajo de eso
queda ligado. El propio techo se calcula en el volumen del primer año de universidad.
\end{remark}

\begin{example}[El rizo]\label{ex:g12:work-and-energy:loop}
Una canica soltada desde la altura $H$ se cuela en un rizo vertical de radio $R$. La
dinámica circular (\cref{ch:g12:newtons-laws}) exige $v_{\text{cima}}^2 \geq gR$
arriba --- el dato que tomaba prestado el problema de la montaña rusa del curso
pasado, ahora nuestro. La conservación desde el reposo en $H$ hasta la cima (altura
$2R$) da $v_{\text{cima}}^2 = 2g(H - 2R)$, luego $H \geq \tfrac52 R$: medio radio por
encima de la cima del rizo, sea cual sea la masa --- y un poco más en el mundo real,
donde se roza.
\end{example}

\begin{omfigure}
\begin{tikzpicture}[font=\small, scale=0.9]
  \draw[thick] (0,3.0) .. controls (0.9,1.1) and (1.6,0) .. (3.1,0) -- (6.9,0);
  \draw[thick] (5.5,1.2) circle (1.2);
  \fill[omDef] (0,3.0) circle (2.2pt);
  \draw[dashed, omExo] (4.0,2.4) -- (7.6,2.4) node[right] {$2R$};
  \draw[dashed, omThm] (-0.4,3.0) -- (7.6,3.0) node[right] {$H_{\min} = \tfrac52 R$};
  \draw[Stealth-Stealth, omThm] (7.35,2.4) -- (7.35,3.0);
  \node[right, omThm] at (7.35,2.7) {$\tfrac{R}{2}$};
  \draw[-Stealth, omMeth, thick] (5.5,2.4) -- (4.5,2.4) node[above right=0pt and 2pt] {$v_{\text{cima}}$};
  \fill (5.5,2.4) circle (1.5pt);
\end{tikzpicture}

{\small El rizo: el punto de suelta queda medio radio por encima de la cima, de modo
que la $E_k$ sobrante mantiene la canica pegada a la pista
($v_{\text{cima}}^2 = gR$).}
\end{omfigure}

\section{Ejercicios}

\begin{exercise}[$\star$]\label{exo:g12:work-and-energy:1}
Un resorte de \omterm{def:g12:oscillators-and-time:spring}{constante elástica} $k = \qty{200}{N/m}$ se estira \qty{5.0}{cm}. Calcula
la \omterm{def:g10:energy-conservation:energy}{energía} almacenada; ¿en qué se convierte si la elongación se duplica? ¿Cuánto
cuestan los \emph{segundos} \qty{5.0}{cm}?
\end{exercise}

\begin{exercise}[$\star$]\label{exo:g12:work-and-energy:2}
Un dron de \qty{0.90}{kg} vuela siguiendo un recorrido caprichoso desde una terraza a
$z = \qty{12}{m}$ hasta un balcón a $z = \qty{30}{m}$. ¿\omterm{def:g11:work-of-force:work}{Trabajo} de su \omterm{def:g10:universal-gravitation:weight}{peso}? ¿Y por la
subida en línea recta? ¿Y en el trayecto de ida y vuelta?
\end{exercise}

\begin{exercise}[$\star$]\label{exo:g12:work-and-energy:3}
¿\omterm{def:g12:work-and-energy:conservative}{Conservativa}, \omterm{def:g12:work-and-energy:conservative}{no conservativa} o sin \omterm{def:g11:work-of-force:work}{trabajo}? Justifícalo con los caminos: (a) el
\omterm{def:g10:universal-gravitation:weight}{peso}; (b) la \omterm{def:g10:inertia:force}{fuerza} de un resorte; (c) el \omterm{def:g10:inertia:inventory}{rozamiento} de deslizamiento; (d) la \omterm{def:g12:newtons-laws:friction}{reacción
normal} sobre una caja que desliza; (e) el \omterm{def:g10:inertia:inventory}{rozamiento} del aire.
\end{exercise}

\begin{exercise}[$\star$]\label{exo:g12:work-and-energy:4}
Un péndulo de \qty{2.0}{m} de longitud se suelta en reposo a $\qty{60}{\degree}$ de la
vertical. Halla su rapidez en el punto más bajo. ¿Por qué la tensión del hilo no entra
nunca en el balance?
\end{exercise}

\begin{exercise}[$\star$]\label{exo:g12:work-and-energy:5}
Una ciclista rueda a \qty{9.0}{m/s} en llano entregando \qty{250}{W}: ¿contra qué
\omterm{def:g10:inertia:force}{fuerza} \omterm{def:g11:work-of-force:sign}{resistente} total pelea? Subiendo a \qty{5.0}{m/s} con la misma \omterm{def:g11:work-of-force:power}{potencia}: ¿qué
\omterm{def:g10:inertia:force}{fuerza} vence?
\end{exercise}

\begin{exercise}[$\star\star$]\label{exo:g12:work-and-energy:6}
El resorte de un lanzador de juguete ($k = \qty{450}{N/m}$) se comprime \qty{8.0}{cm}
detrás de un dardo de \qty{20}{g}. Halla la \omterm{def:g10:energy-conservation:energy}{energía} almacenada, la rapidez de
lanzamiento y la altura alcanzada si se dispara verticalmente (sin aire).
\end{exercise}

\begin{exercise}[$\star\star$]\label{exo:g12:work-and-energy:7}
Estirando una goma elástica, un estudiante mide $F = 0$, $3.0$, $7.0$, $12.0$,
\qty{18.0}{N} en $x = 0$, $2.0$, $4.0$, $6.0$, \qty{8.0}{cm}. Estima la \omterm{def:g10:energy-conservation:energy}{energía}
almacenada por áreas de trapecios; compárala con $\tfrac12 kx^2$, tomando $k = F/x$ en
el extremo. ¿Por qué difieren?
\end{exercise}

\begin{exercise}[$\star\star$]\label{exo:g12:work-and-energy:8}
Un trineísta de \qty{55}{kg} parte del reposo, baja \qty{8.0}{m} a lo largo de una
pista curva de \qty{25}{m} y llega a \qty{9.0}{m/s}. ¿\omterm{def:g11:mechanical-energy:em}{Energía mecánica} perdida?
¿\omterm{def:g10:inertia:force}{Fuerza} media de \omterm{def:g10:inertia:inventory}{rozamiento}? ¿Por qué es irrelevante la forma exacta de la pista?
\end{exercise}

\begin{exercise}[$\star\star$]\label{exo:g12:work-and-energy:9}
Una pista de canicas tiene un rizo de \qty{20}{cm} de radio. Halla la altura mínima de
suelta (desde el reposo, sin \omterm{def:g10:inertia:inventory}{rozamiento}). Soltada desde \qty{60}{cm}: ¿rapidez en la
cima del rizo y cociente $v_{\text{cima}}^2/(gR)$? ¿Por qué una pista real debe
lanzarse desde aún más arriba?
\end{exercise}

\begin{exercise}[$\star\star$]\label{exo:g12:work-and-energy:10}
Una cuenta de \qty{100}{g} desliza sin \omterm{def:g10:inertia:inventory}{rozamiento} por un eje, con
$E_p(x) = 8.0\,x^2$ (en \omterm{def:g11:work-of-force:work}{julios}, con $x$ en \omterm{def:g10:orders-of-magnitude:unit}{metros}) y $E_m = \qty{2.0}{J}$. ¿\omterm{def:g12:work-and-energy:diagram}{Puntos de
retorno}? ¿Rapidez máxima? ¿Qué tipo de movimiento es este?
\end{exercise}

\begin{exercise}[$\star\star$]\label{exo:g12:work-and-energy:11}
Un coche de \qty{1200}{kg} a \qty{2.0}{m/s} es detenido por el resorte de un parachoques
de \omterm{def:g12:oscillators-and-time:spring}{constante elástica} $k = \qty{6.0e5}{N/m}$. Halla la compresión máxima. A
\qty{4.0}{m/s}: ¿por qué no es cuatro veces mayor, si la \omterm{def:g10:energy-conservation:energy}{energía} se cuadruplica?
\end{exercise}

\begin{exercise}[$\star\star\star$]\label{exo:g12:work-and-energy:12}
Un carrito de \qty{0.50}{kg} rueda por una pista cuya $E_p(x)$ tiene una cima de
\qty{6.0}{J} ($x = 0$) y un valle de \qty{2.0}{J} ($x = \qty{1.5}{m}$). (a) Si llega
por la izquierda con $E_m = \qty{5.0}{J}$: ¿pasa? ¿Qué hace en su lugar? (b) Con
$E_m = \qty{7.0}{J}$: ¿rapideces en la cima y en el valle? (c) ¿Equilibrios?
\end{exercise}

\begin{exercise}[$\star\star\star$]\label{exo:g12:work-and-energy:13}
Admitiendo que subir una masa $m$ desde el suelo hasta una distancia arbitrariamente
grande cuesta $mgR_E$, con $R_E = \qty{6.37e6}{m}$ (el área finita bajo la gráfica de
la gravedad menguante), calcula la \emph{\omterm{rem:g12:work-and-energy:escape}{velocidad de escape}} de la Tierra. ¿Por qué
es la misma para una sonda que para un guijarro? ¿Y un lanzamiento justo por debajo de
ella?
\end{exercise}

\begin{exercise}[$\star\star\star$]\label{exo:g12:work-and-energy:14}
Un móvil de \qty{250}{g} oscila sobre un resorte ($k = \qty{40}{N/m}$) con \omterm{def:g12:oscillators-and-time:oscillator}{amplitud}
\qty{6.0}{cm} (\cref{ch:g12:oscillators-and-time}). ¿\omterm{def:g10:energy-conservation:energy}{Energía} total? ¿Rapidez máxima?
¿Posiciones donde $E_k = E_p$?
\end{exercise}

\begin{exercise}[$\star\star\star$]\label{exo:g12:work-and-energy:15}
Un saltador de pértiga esprinta a \qty{9.5}{m/s}. Estima la altura que gana su \omterm{def:g11:forces-and-motion:com}{centro
de masas} si la pértiga convierte toda su \omterm{def:g11:mechanical-energy:kinetic}{energía cinética}, y el listón superado (con
el \omterm{def:g11:forces-and-motion:com}{centro de masas} partiendo de \qty{1.0}{m}). El récord anda por \qty{6.3}{m}: ¿qué
capta la estimación y qué deja fuera?
\end{exercise}

\section{Problema: El trampolín de saltos de esquí}

\begin{problem}[{Problema de fin de semana --- el trampolín de esquí: una rampa
auditada con un radar de velocidad, un vuelo en caída libre, la suave geometría de una
recepción empinada y la torre que un proyectista tiene que comprar cuando el
rozamiento pasa factura}]
\label{pb:g12:work-and-energy:1}
Un saltador de esquí olímpico ($m = \qty{70}{kg}$, esquís incluidos) parte del reposo
desde una puerta situada $H = \qty{45}{m}$ por encima de la mesa de despegue, desliza
$L = \qty{90}{m}$ de rampa curva y deja la mesa horizontalmente; allí un \omterm{met:g10:signals-and-waves:echo}{radar} de
velocidad marca $v_0 = \qty{26}{m/s}$. El punto de recepción queda $h = \qty{40}{m}$
más abajo, sobre una pendiente de $\qty{35}{\degree}$. Las partes II y III desprecian
la resistencia del aire.

\medskip\noindent\textbf{Parte I --- La rampa.}
\begin{enumerate}
  \item Predice la rapidez de despegue sin \omterm{def:g10:inertia:inventory}{rozamiento}; ¿por qué es irrelevante la
        curva de la rampa?
  \item Con la referencia en la mesa, calcula $E_m$ en la puerta y en el despegue.
  \item Calcula $\Delta E_m$; ¿qué teorema nombra a los culpables, y quiénes son?
  \item Deduce la \omterm{def:g10:inertia:force}{fuerza} \omterm{def:g11:work-of-force:sign}{resistente} media (nieve más aire) a lo largo de la rampa.
  \item ¿Qué fracción de $E_m$ se pierde? ¿Por qué encerar los esquís y agacharse?
\end{enumerate}

\medskip\noindent\textbf{Parte II --- El vuelo.}
\begin{enumerate}[resume]
  \item En el vuelo solo actúa el \omterm{def:g10:universal-gravitation:weight}{peso}: ¿de qué movimiento del
        \cref{ch:g12:projectile-motion} se trata? ¿Componentes de la velocidad en la
        recepción?
  \item La rapidez de recepción, dos veces: a partir de las componentes y a partir de
        la conservación de $E_m$. ¿En \unit{km/h}?
  \item Calcula el \omterm{def:g12:projectile-motion:range}{tiempo de vuelo} y la distancia horizontal recorrida.
  \item Los saltadores reales cabalgan el aire como un ala y vuelan mucho más lejos:
        ¿puede el aire \emph{aumentar} la rapidez de recepción? Argumenta con el
        \omterm{thm:g12:work-and-energy:emtheorem}{teorema de la energía mecánica}.
\end{enumerate}

\medskip\noindent\textbf{Parte III --- La recepción.}
\begin{enumerate}[resume]
  \item Calcula el ángulo de la velocidad de recepción por debajo de la horizontal.
  \item Descomponla en las componentes paralela y perpendicular a la pendiente.
  \item Las piernas solo absorben la parte perpendicular: calcula esa \omterm{def:g11:mechanical-energy:kinetic}{energía
        cinética}; compárala con la total.
  \item Un suelo horizontal lo absorbería todo: ¿a qué caída vertical equivale cada
        recepción ($h_{\text{eq}} = v_\perp^2/2g$, resp.\ $v^2/2g$)?
  \item En una frase: ¿por qué las pistas de recepción son empinadas y curvas?
\end{enumerate}

\medskip\noindent\textbf{Parte IV --- La torre del proyectista.}
Un trampolín más pequeño necesita despegar solo a $v_0' = \qty{24}{m/s}$; su rampa
recta desciende a $\qty{35}{\degree}$, así que una altura de puerta $H'$ significa una
longitud de pista $H'/\sin\qty{35}{\degree}$; el mismo \omterm{def:g10:inertia:inventory}{rozamiento} de \qty{80}{N}.
\begin{enumerate}[resume]
  \item Calcula la altura de puerta que construiría un proyectista sin \omterm{def:g10:inertia:inventory}{rozamiento}.
  \item Escribe el \omterm{thm:g12:work-and-energy:emtheorem}{teorema de la energía mecánica}, de la puerta a la mesa, para $H'$.
  \item Despeja $H'$.
  \item ¿En qué porcentaje elevó el \omterm{def:g10:inertia:inventory}{rozamiento} la torre?
  \item La masa ya no se cancela: recalcula $H'$ para un saltador de \qty{60}{kg}.
        ¿Quién necesita la torre más alta, y por qué?
  \item El remate: la torre anunciada y la parte de ella que se lleva el \omterm{def:g10:inertia:inventory}{rozamiento}.
\end{enumerate}
\end{problem}
